\documentclass[reqno]{amsart}
\usepackage[T1]{fontenc}
\usepackage{amsmath,amssymb,amsthm}
\usepackage{xr-hyper}
\usepackage[hidelinks]{hyperref}
\externaldocument[ova-]{20261125T001000--zeros-to-ramp-sums__workbook}
\externaldocument[ovb-]{20261125T001100--inert-primes-and-estimate-b__workbook}
\externaldocument[ovc-]{20261124T223100--complex-multiplication-and-inert-reduction__workbook}
\externaldocument[ovd-]{20261124T214500--supersingular-theta-divisibility__workbook}
\externaldocument[ove-]{20261126T030200--the-bound-is-attained__workbook}
\externaldocument[ovf-]{20261125T007000--stepanov-without-multiplicity__workbook}
\externaldocument[ovg-]{20261126T030300--the-data__workbook}
\externaldocument[ovh-]{20261125T005000--galois-fourier-and-averaging__workbook}
\externaldocument[ovi-]{20261126T031100--theorem-c-and-the-main-inequality__workbook}
\externaldocument[ovj-]{20261125T003000--the-chain-and-the-square-root__workbook}
\externaldocument[ovk-]{20261126T040100--spaces-forms-and-twists__workbook}
\externaldocument[ovl-]{20261126T030600--theorem-a-short-forms__workbook}
\externaldocument[ovm-]{20261125T010000--rank-lemma-and-zero-estimate__workbook}
\externaldocument[ovn-]{20261126T020210--bessel-poincare-series-and-their-gram-operators__workbook}
\externaldocument[ovo-]{20261126T030100--a-twist-for-each-prime__workbook}
\externaldocument[ovp-]{20261126T040200--theorem-c__workbook}
\externaldocument[ovq-]{20261126T030700--the-determinant-of-a-chain__workbook}
\externaldocument[ovr-]{20261126T030800--the-cost-of-one-lift__workbook}
\externaldocument[ovs-]{20261126T031000--twists-with-few-congruences__workbook}
\externaldocument[ovt-]{20261126T052000--which-terms-survive-characters-minors-and-the-gain-at-the-level__workbook}

% The chapter has a name and its own section and problem numbers.
\newcounter{chap}
\renewcommand{\thechap}{Overview}
\theoremstyle{definition}
\newtheorem{problem}{Problem}[section]

\newcommand{\rem}{\mathrm{rem}}
\newcommand{\Q}{\mathbb{Q}}
\setlength{\emergencystretch}{3em}

\title{Overview}
\author{}
\date{November 27, 2026}

\begin{document}
\maketitle
\refstepcounter{chap}\label{chap}

This chapter is optional.  It outlines the proof in a few pages and
points to the chapters where each step is carried out or stated.
Nothing later depends on it.

\section{The proof in one paragraph}\label{sec:proof}

An upper bound for one sum over primes excludes the zeros of
Dirichlet \(L\)-functions in the critical strip off the line
\(\Re s = 1/2\).  The bound comes from comparing the size of a
determinant with its divisibility, as in Chebyshev's estimates for
primes.  The determinant is a Gram determinant of automorphic forms
after a chain of \(m\) theta lifts, where \(m \approx \sqrt{x}\) and
\(x\) is the cutoff in the prime sum.  Its entries involve
derivatives of theta functions with complex multiplication by
\(K = \Q(\sqrt{-d})\).  These derivatives are highly divisible at
the primes that do not split in \(K\), and this is how primes
enter.  The remainder of the prime sum occurs once for each lift,
\(m\) times in all, and the proof bounds the total of the errors for
the whole chain.  Dividing by \(m\) gives the exponent \(-1/2\).

\section{From zeros to a sum over primes}\label{sec:zeros}

Let \(q_0\) be a multiple of \(18480\), let \(d\) be a prime with
\(d \equiv q_0/2 - 1 \pmod{q_0}\), and let \(\chi_K\) be the
quadratic character of \(K = \Q(\sqrt{-d})\).  It is completely
multiplicative, and at a prime \(p\) its value is \(1\), \(-1\) or
\(0\) according as \(p\) splits, is inert, or ramifies in \(K\).
For \(y \ge 1\) let
\[
  R_{q_0}(y) :=
  \sum_{\substack{n \le y \\ n \equiv 1 \, (q_0)}}
  \Lambda(n) (1 - \chi_K(n)) \Bigl( \frac{1}{n} - \frac{1}{y} \Bigr)
  - \frac{\log y}{\varphi(q_0)} - c ,
\]
where \(\Lambda\) is the von Mangoldt function, \(\varphi\) is Euler's
totient, and \(c = c_{d,q_0}\) is the constant for which
\(R_{q_0}(y) \to 0\).
The target is the one-sided bound
\begin{equation}\label{eq:target}
  R_{q_0}(y) \le y^{-1/2 + \varepsilon}
  \quad \text{for all large } y,
\end{equation}
for every \(\varepsilon > 0\) and every such \(d\).

Why does an upper bound suffice?  The Mellin transform of
\(R_{q_0}\) can have poles at the zeros of the \(L(s, \psi)\), with \(\psi\)
modulo \(q_0\), and at the zeros of their twists \(L(s, \psi\chi_K)\),
whose multiplicities enter with opposite signs.  Given a zero \(\rho\) with
\(1/2 < \Re \rho < 1\), one chooses \(d\) so that no twist vanishes in a
region containing \(\rho\) and the real interval \([\Re \rho, 1]\).  If \eqref{eq:target} holds, the function
\(y^{-1/2+\varepsilon} - R_{q_0}(y)\) is eventually nonnegative.
For \(\varepsilon\) small enough, its transform still has the pole
at \(\rho\).  Landau's theorem forces a real singularity at or to
the right of \(\Re \rho\), where the only possible poles of this
transform have negative residue.  A negative residue contradicts
eventual nonnegativity.  For a character of modulus \(q\) one takes
\(q_0 = 18480q\).
See Chapter~\ref{ova-chap} (\emph{From zeros to a one-sided bound}) and
Chapter~\ref{ovb-chap} (\emph{Inert primes and the choice of \(d\)}).

The weight \(1 - \chi_K(n)\) vanishes at the powers of the primes
that split in \(K\).  For the other primes, the sum of the weights
of the powers \(p^h\), each divided by \(p^h\) as in
\(R_{q_0}\), is as follows.

\begin{problem}\label{pr:weight}
  Let \(p\) be a prime.  The sum of \((1 - \chi_K(p^h))/p^h\) over
  \(h \ge 1\) is \(2p/(p^2 - 1)\) if \(p\) is inert in \(K\), and
  \(1/(p - 1)\) if \(p\) ramifies.
\end{problem}

\section{Size against divisibility}\label{sec:divisibility}

A nonzero integer divisible by a positive integer \(D\) has absolute
value at least \(D\).  Chebyshev applied this to \(\binom{2n}{n}\),
which is at most \(4^n\) and divisible by every prime in \((n, 2n]\).

\begin{problem}\label{pr:chebyshev}
  For every integer \(n \ge 1\), \(\sum_{n < p \le 2n} \log p \le n \log 4\).
\end{problem}

The proof applies the same principle to derivatives of theta
functions.  The model is \(u!\).  Its logarithm is
\(u \log u - u + O(\log u)\), and Legendre's formula
\(\log u! = \sum_{n \le u} \Lambda(n) \lfloor u/n \rfloor\) gives its
divisibility: the prime power \(n\) contributes about \(u/n\).

The theta functions of the proof have complex multiplication by
\(K\) and variables indexed by the infinite places \(\tau\) of a CM
field \(E \supseteq K\).  Each place has a \emph{weight}
\(\kappa(\tau)\), a nonnegative integer that sets the degree of
differentiation in the variables at \(\tau\), and at some places there
are extra derivatives of a large order \(u\).  The sizes of such
derivatives are governed by factorials.  Their divisibility comes
from Frobenius at the primes that do not split in \(K\).  A power
\(n\) of such a prime relates a place \(\tau\) to a second place
\([-n]\tau\), its \emph{partner}, and the bound for the divisibility
of the derivatives of order \(u\) at \(\tau\) is, on a logarithmic
scale and up to smaller losses,
\[
  \sum_{n} \Lambda(n) (1 - \chi_K(n))
  \Bigl( \frac{u}{n} - \kappa([-n]\tau) \Bigr)_+ ,
\]
where \(t_+ := \max(t, 0)\).  The term \(u/n\) is the gain at \(n\), as for \(u!\).  Frobenius
couples \(n\) degrees at \(\tau\) with one degree at the partner,
which has \(\kappa([-n]\tau)\) degrees; the positive part is the
excess left uncoupled.  For orders near the weights the bound is
attained, up to a common scalar.  See
Chapter~\ref{ovc-chap} (\emph{Complex multiplication and reduction at an inert prime}),
Chapter~\ref{ovd-chap} (\emph{The invariant derivation at a supersingular prime}) and
Chapter~\ref{ove-chap} (\emph{The bound is attained}).

The subtraction produces a cutoff.  If the order is \(a := xb\) and
the partner has weight \(b\), only \(n \le x\) contribute, each
with \(a\) times the factor \(1/n - 1/x\) of the sum defining
\(R_{q_0}(x)\):

\begin{problem}\label{pr:ramp}
  Let \(x, b, n > 0\) and \(a = xb\).  Then
  \((a/n - b)_+ = a (1/n - 1/x)_+\).
\end{problem}

If every partner had weight \(b\), the sum would be
\(a (\log x + O(1))\), by Mertens' estimate
\(\sum_{n \le y} \Lambda(n)(1 - \chi_K(n))/n = \log y + O(1)\); more
precisely, it would be an explicit main term plus \(a\) times a
remainder of the kind bounded in \eqref{eq:target}.  The sizes are
arranged to cancel the main term; this is the role of one of the
levels of the forms.  ``Divisibility is at most size'' then bounds
the remainder from above.
Chapter~\ref{ovf-chap} (\emph{Stepanov without multiplicity: a determinant against Siegel zeros})
gives models for this comparison with one determinant.

\section{The chain}\label{sec:chain}

To make the remainder occur \(m\) times, the proof composes \(m\)
lifts.  A theta lift is a map between spaces of automorphic forms
on two unitary groups whose kernel is a theta function.  The lifts
of the proof pass through unitary groups of increasing rank, the
last of them definite, and the proof takes Gram determinants of
the images of a family of forms.  Each lift contributes the
divisibility above.  The outcome is
\begin{equation}\label{eq:main}
  m \cdot \rem \le f a (\log b)^5 ,
\end{equation}
with the following notation.  The scales are \(x < x_1 < b\), and
\(a = xb\).  The field \(E\) is cyclotomic over \(K\), with \(f\)
infinite places.  The weight is \(b\) on a set \(\Phi_b\) of places
and \(a/x_1\) on the others, and the extra derivatives, of order
\(a\), are at a set \(\Phi_a\) of places.  For each
\(n \equiv 1\) modulo \(q_0\), exactly \(\lvert \Phi_b \rvert\) of
the places \(\tau \in \Phi_a\) have their partner in \(\Phi_b\),
with cutoff \(x\); for all other pairs \((n, \tau)\) the cutoff is
\(x_1\).  The quantity \(\rem\) is the sum over \(\tau \in \Phi_a\)
of the divisibility sums of order \(a\), less their main term, and
the count gives
\begin{equation}\label{eq:rem}
  \frac{\rem}{a} = \lvert \Phi_b \rvert \bigl( R_{q_0}(x) - R_{q_0}(x_1) \bigr)
  + \lvert \Phi_a \rvert R_1(x_1) ,
\end{equation}
where \(R_1\) is the remainder with modulus \(1\).  The prime number
theorem in progressions bounds the remainders at \(x_1\), and
\(\log b\) is at most a small power of \(x\), so \eqref{eq:main}
bounds \(R_{q_0}(x)\) by a small power of \(x\) divided by \(m\).
See Chapter~\ref{ovg-chap} (\emph{The data}),
Chapter~\ref{ovh-chap} (\emph{The places, the weights and the remainder}) and
Chapter~\ref{ovi-chap} (\emph{The main inequality}).

The costs of the chain limit \(m\): two of them, relative to what
is available, are of order \(m^2/x\) up to powers of \(\log b\), and
the proof takes \(\sqrt{x} \le m < 2\sqrt{x}\).  See
Chapter~\ref{ovj-chap} (\emph{The chain and the square root}), which
starts from the following model.

\begin{problem}\label{pr:sqrt}
  For \(A, x > 0\), the minimum of \(A (1/m + m/x)\) over real
  \(m > 0\) is \(2A/\sqrt{x}\), attained at \(m = \sqrt{x}\).
\end{problem}

\section{Three theorems}\label{sec:theorems}

The inequality \eqref{eq:main} follows from three statements.  They
concern a space \(S_{\mathrm{ram}}\) of forms on a definite unitary
group, the source; a space \(T\) of dimension \(d_T\) of forms on
the first group of the chain, the target; maps
\(S_{\mathrm{ram}} \to T\), the translates of one map defined by
Poincar\'e series; and the chain on \(T\).  The forms have integral
structures, and the chain can be twisted by a character; the
trivial twist is the chain itself.  These objects are defined in
Chapter~\ref{ovk-chap} (\emph{Spaces, forms and twists}).

\emph{Theorem A: short forms.}  The integral forms in
\(S_{\mathrm{ram}}\) whose images in \(T\) have short lifts span
\(S_{\mathrm{ram}}\), and families of images of integral forms have
bounded index: a change of columns of bounded determinant makes
them integral.  Here ``short'' means that the logarithm of the
squared length of the lift is at most \(H - m \cdot \rem + F\),
where \(H \ge 0\) is the main term and \(F := f a (\log b)^4\); the
divisibility of the theta derivatives has been divided out.  The
chain on \(T\) starts on an indefinite group.  A seesaw identity
expresses its composition with the maps \(S_{\mathrm{ram}} \to T\),
up to a scalar, through a chain of lifts between definite unitary
groups followed by a restriction.  The kernels of the definite
lifts are derivatives of CM theta functions, whose size and
divisibility can be computed.  See Chapter~\ref{ovl-chap} (\emph{Theorem A: short forms}).

\emph{Theorem B: enough forms.}  The images of \(S_{\mathrm{ram}}\)
under a few translates span a subspace of \(T\) of codimension at
most \(4 d_T/b^2\).  The proof is a second moment estimate for the
Poincar\'e series: on average the translates are nearly
orthogonal, so, as long as much of \(T\) is missing, some translate
reduces the missing dimension by a fixed proportion.  The proof
also shows that \(d_T > 0\) and bounds \(d_T\) from above; the
bound uses the equal multiplicities of Part~IX.  See
Chapter~\ref{ovm-chap} (\emph{Few translates span}) and
Chapter~\ref{ovn-chap} (\emph{Theorem~B: enough forms}).

\emph{Theorem C: the lower bound.}  For an independent family in
\(T\) spanning a subspace of codimension at most \(4 d_T/b^2\), let
\(c\) bound the logarithm of the Gram determinant of the lifts at
every conjugate of the coordinates, and let \(\delta\) bound the
logarithm of the index.  Then \(c + 2\delta \ge d_T (H - 3F)\).
This replaces ``a nonzero integer has
absolute value at least \(1\)''.  Integrality is not preserved
exactly along the chain: projections onto images of lifts have
denominators, at primes where forms are congruent.  Sizes are
therefore read at the trivial twist, and integrality, for each
prime and each embedding of the coefficients above it, at a twist
chosen to have few congruences.  The twist may vary because, for a
basis of \(T\), the Gram determinant at the trivial twist is the
Gram determinant at a twist \(\vartheta\) times one scalar, the
ratio of the determinants of the two chains.  This ratio is the
product of the ratios for the single lifts, up to a power of the conductor of the twist and ratios of Hecke
\(L\)-values, and the ratio for one lift is bounded using forms for
which the divisibility is attained.  These forms also account for
\(H\).  Besides the derivatives of large order, each lift has
derivatives of the orders given by the weights, and their size
exceeds their guaranteed divisibility.  A pays this excess, \(H\)
in all, in the length; the forms that attain the divisibility
return it in the lower bound.  Only the remainder from the
derivatives of large order is left.  See
Chapter~\ref{ovo-chap} (\emph{A twist for each prime}),
Chapter~\ref{ovp-chap} (\emph{Theorem C}),
Chapter~\ref{ovq-chap} (\emph{The determinant of a chain}),
Chapter~\ref{ovr-chap} (\emph{The cost of one lift}) and
Chapter~\ref{ovs-chap} (\emph{Twists with few congruences}).

By B and A there are \(n \ge d_T/2\) independent vectors in \(T\),
images of integral forms with short lifts, spanning a subspace of
codimension at most \(4 d_T/b^2\).  Let
\(F' := f a (\log b)^3 \log(1 + 2b^2)\).  By Hadamard's inequality
(Problem~\ref{pr:hadamard}), A gives the bound
\(c := n (H - m \cdot \rem + F)\) for the logarithm of the Gram
determinant of the lifts, at every conjugate.  A bounds the index for each translate
and B the number of translates; together they bound the logarithm
of the index by \(\delta := 10 d_T F'\).  The index enters twice,
once for each side of the Gram pairing
(Problem~\ref{pr:columns}).  Thus
\begin{equation}\label{eq:compare}
  \begin{aligned}
  d_T (H - 3F)
  &\overset{\mathrm{C}}{\le} c + 2\delta \\
  &\overset{\mathrm{A, B}}{\le}
  d_T (H + F) - n m \, \rem + 20 \, d_T F' .
  \end{aligned}
\end{equation}
Subtracting \(d_T H\) leaves \eqref{eq:main}
(Problem~\ref{pr:compare}), and \eqref{eq:main} gives the target at
integers (Problem~\ref{pr:final}).  The passage from integers to
real \(y\) is in Chapter~\ref{ovi-chap} (\emph{The main inequality}).

\begin{problem}[Hadamard's inequality]\label{pr:hadamard}
  The Gram determinant \(\det(\langle v_i, v_j \rangle)\) of vectors
  \(v_1, \dots, v_n\) in a positive definite Hermitian space is at
  most \(\prod_j \lVert v_j \rVert^2\).
\end{problem}

\begin{problem}\label{pr:columns}
  For an \(n \times n\) complex matrix \(P\), replacing the vectors
  \(v_j\) by \(\sum_i P_{ij} v_i\) multiplies their Gram determinant
  by \(\lvert \det P \rvert^2\).
\end{problem}

\begin{problem}\label{pr:compare}
  Let \(d_T > 0\), \(d_T/2 \le n \le d_T\) and \(F, F' \ge 0\).  Then
  \eqref{eq:compare} gives \(m \cdot \rem \le 8F + 40F'\).  If
  \(F = f a (\log b)^4\), \(F' \le 3F\) and \(\log b \ge 128\), this
  gives \eqref{eq:main}.
\end{problem}

\begin{problem}\label{pr:final}
  Suppose \eqref{eq:main} and \eqref{eq:rem}, with
  \(f \le \varphi(q_0) \lvert \Phi_b \rvert \log b\),
  \(\lvert \Phi_a \rvert \le \lvert \Phi_b \rvert \log b\),
  \(m \ge \sqrt{x}\) and \(\log b \le x^{\eta} + 1\), and suppose
  that \(\lvert R_{q_0}(x_1) \rvert\) and
  \(\lvert R_1(x_1) \rvert \log b\) are at most \(1/x\).  Then
  \[
    R_{q_0}(x) \le \varphi(q_0) (x^{\eta} + 1)^6 x^{-1/2} + 2/x .
  \]
  If this holds for every large integer \(x\), with
  \(6\eta < \varepsilon\), then \eqref{eq:target} holds at the large
  integers.
\end{problem}

\section{How the Parts fit together}\label{sec:parts}

The table groups the chapter arrows in the \href{flowchart.html}{flowchart} by Part.  Part VI specifies the data, spaces and forms, and combines the three main theorems to obtain the main inequality.

\begin{center}
\begin{tabular}{@{}l p{0.47\linewidth} p{0.36\linewidth}@{}}
  \textbf{Part} & \textbf{What it does} & \textbf{Needs} \\
  I & The target & --- \\
  II & The method & I \\
  III & Divisibility at inert primes & I, II \\
  IV & Sizes, periods and counts & II, III \\
  V & Denominators and congruences & II, III \\
  VI & Data, spaces and the main inequality & I, II, VII, VIII, IX \\
  VII & Theorem B: enough forms & II, IV, IX \\
  VIII & Theorem A: short forms & II, III, IV \\
  IX & Theorem C: the lower bound & I, II, III, IV, V \\
\end{tabular}
\end{center}

\section{Where the difficulty is}\label{sec:difficulty}

The comparison and the passage from \eqref{eq:main} to
\eqref{eq:target} are elementary.  The difficulty is in the three
theorems, in particular in the following.
\begin{enumerate}
\item In C, one twist at each prime must bound the denominators of
  the projections and the ratios of Hecke \(L\)-values for all
  \(m\) lifts.
\item In B, a bound for the second moment of the right order of
  magnitude is not enough: it must nearly equal the bound that the
  dimension of \(T\) imposes.
\item In A and C, the divisibility of the derivatives must be
  bounded in several variables at once, and there must be forms
  for which the bound is attained.
\item In C, the identities between lifts that factor the
  determinant of the chain.
\end{enumerate}

\section{What is taken on trust}\label{sec:trust}

The sheets prove the deductions (the theorem from Theorems A, B and C, and
  each of these from its inputs) and, for most inputs, a simplest instance
  or a model.  The following are stated from the Lean proof and not proved;
  classical results are cited in each sheet.
\begin{description}
\item[Part I] the prime number theorem in progressions with every power of
  the logarithm saved (classical).
\item[Part VI] the existence of the data; the finite dimension of the spaces
  of forms under the gap hypothesis.
\item[Part VII] the block identity; the lower bound for the Bessel period at
  a real place and the phase and cone estimates used for it; the absolute
  convergence of the Bessel average for the data, the decay majorant and the
  cut with its tail bound; the two inequalities between the identity and
  main amounts; which terms survive, and the size bound for a cut block in
  terms of its near representatives, the gain and count of pairs, and the
  tail of the cut, with the stability and summation properties of the
  pieces; for the open cell, the normal form over an open reduction, the
  bound for the sum of the head weights, the comparison of constants, and
  the local bounds at finite and infinite places for pairs of translates of the \emph{Bessel tuple} \((v_1,\dots,v_\varrho,z)\) of Definition~\ref{ovt-def:frame}; the
  facts about the Bessel subgroup and its representation that the unfolding
  and the projection to good forms use (it intertwines the translates; its
  trace; the conjugacy average).
\item[Part VIII] the steps of the seesaw identity; the expansion of a chain
  datum into a weak one, with the integrality and the additivity of
  exponents along the chain; the four facts about the CM theta function and
  its one-curve expansion at a non-split prime; the theta-product
  factorization for a standard datum and jet integrality at primes dividing
  the cyclotomic conductor; the passage from derivatives to one lift; the
  count with charges, and three facts of the general step of the count; the
  theorem on elementary divisors, with, in its proof, the weighted
  Nullstellensatz, the general root lemma, the count of forms with bounded
  defects and the average; the steps at the exceptional primes and at the
  level prime; the estimate for the restricted chain behind the size (the
  bound for one kernel on its conjugates, the class count at congruence
  level, the branching norm); the integrality of the Weyl pairing; the
  commutation of the lower chain with translations at the chosen places.
\item[Part IX] the transfer of periods along a turn and the spanning forms
  with short lifts, with the estimates on the last path (behind the
  auxiliary forms); the inputs of the Gram identity of one lift; the
  comparison of trace formulas behind the equality of Hecke traces; the
  exchange identity (only its scalar is computed here) and three statements
  on a triangle of lifts (its realization, the inner product formula on its
  first edges, the comparison at the last level); the Hecke relation of the
  lifts and the projection estimate (behind the compressions); the period
  identity and the bounds for the gluing maps, that a form is detected by
  its common periods, and that false lifts have no periods (behind the
  comparison operator); the norm of the lift of a character; the bounds on
  six and twelve copies; the pairing of test forms; Galois representations
  of eigenforms on definite unitary groups and their local properties above
  \(\ell\) and at tame level, stated with references in Part V; the bound
  for the loss at one place above \(\ell\); the bound for the eigenvalues of
  one clearing operator behind the denominators of projections, with the
  first-order lemma and, at the auxiliary places, the triangle of primes and
  the first-order operators; the existence of main summands; the denominator
  clause of regularity; the inputs to the twisting theorem (a finite
  presentation of the cocycles, the release of the ties at the first prime,
  and the auxiliary primes and local bounds behind the tied primes; the
  general measure for Hecke ratios, including the determinant period at the
  edge value, and its Stickelberger systems; existence of lattices, the
  crystalline cocycle rank formula and transport of crystalline relations;
  and the return from Galois characters to tame characters).
\end{description}

Classical background on Galois cohomology and crystalline lattices is stated
  with references in Part V.  Part VII gives references for discrete series
  of real unitary groups.

\section{Notation}\label{sec:notation}

Each sheet defines the notation it uses; the following conventions are
  common to all sheets.
\begin{description}
\item[\(q_0,d,K,\chi_K\)] \(q_0\) and \(d\) are the parameters and \(K =
  \mathbb{Q}(\sqrt{-d})\), with character \(\chi_K\).
\item[\(E,f,\Phi_a,\Phi_b,\kappa(\tau)\)] \(E\) is the CM field of the data,
  with \(f\) infinite places, among them the \(a\)-places \(\Phi_a\) and the
  \(b\)-places \(\Phi_b\), and \(\kappa(\tau)\) is the weight of a place.
\item[\(x,x_1,b,a\)] The scales are \(x < x_1 < b\), with \(a = xb\).
\item[\(m,r,k,t\)] \(m\) is the number of lifts, \(r\) and \(k = r + 1 - m\)
  are the ranks of the target and of the source, and \(t = 5r\): the ranks
  of the spaces \(Y_j\) increase by \(t\), those of the \(X_j\) by \(t +
  1\).
\item[\(M,N\)] \(M = p_0p_1\) and the prime \(N\) are the levels.
\item[\(h_{\mathrm{main}},\rem,U\)] The main amount is \(h_{\mathrm{main}} =
  mk\,h_{\mathrm{dir}}\), the remainder is \(\rem\), and \(U = fa\log^3 b\)
  is the unit of the errors.
\item[\(\vartheta,\theta\)] A twist is \(\vartheta\); a family of twists
  \(\theta = (\theta_\iota)\) has one member for each isomorphism
  \(\iota\colon \mathbb{C} \simeq \bar{\mathbb{Q}}_\ell\).
\item[\(S,S_{\mathrm{ram}},T,d_T\)] \(S \supseteq S_{\mathrm{ram}}\) are the
  source forms and the ramified ones, \(T\) the target forms, of dimension
  \(d_T\).
\item[\(\beta,\operatorname{val}\)] The diagram \(\beta\) gives the weights
  of the forms; \(\operatorname{val}\) is the valuation with
  \(\operatorname{val}(p) = 1\) after a ring map to \(\mathbb{C}_p\).
\item[\(R_{d,q}\)] The remainder \(R_{d,q}\) of the prime sum is as in
  Chapter~\ref{ova-chap} (\emph{From zeros
  to a one-sided bound}).
\end{description}
The numbers are defined in
  Chapter~\ref{ovg-chap} (\emph{The data}), the
  spaces, forms and twists in
  Chapter~\ref{ovk-chap} (\emph{Spaces,
  forms and twists}); a chapter that uses one of these letters otherwise
  says so.

\section*{Hints}
\begin{description}
\item[\ref{pr:weight}] \(\chi_K(p^h) = \chi_K(p)^h\).  Sum a geometric
  series over the odd \(h\), or over all \(h\).
\item[\ref{pr:chebyshev}] \(\binom{2n}{n}\) is one of the \(2n + 1\)
  terms of \((1 + 1)^{2n}\), and each prime in \((n, 2n]\) divides
  \((2n)!\) but not \(n!^2\).
\item[\ref{pr:ramp}] \(b = a/x\).
\item[\ref{pr:sqrt}] Apply the inequality of the arithmetic and
  geometric means to \(1/m\) and \(m/x\).
\item[\ref{pr:hadamard}] For independent vectors, apply
  Gram--Schmidt.  It does not change the determinant and does not
  increase lengths.
\item[\ref{pr:columns}] With the pairing conjugate-linear in the first
  variable, the new Gram matrix is \(P^* G P\).
\item[\ref{pr:compare}] Subtract \(d_T H\), then divide by \(n\).
\item[\ref{pr:final}] Divide \eqref{eq:main} by
  \(m a \lvert \Phi_b \rvert\) and solve \eqref{eq:rem} for
  \(R_{q_0}(x)\).
\end{description}

\end{document}
